% !TeX root = ../../main.tex
\section{Der Begriff Smart-Home}
Für den Begriff Smart Home existiert keine einheitliche und allgemein anerkannte Definition.  
In der Fachliteratur werden unterschiedliche Begriffe wie Connected Home, Smart Living oder intelligentes Wohnen teilweise synonym verwendet.  
Gemeinsam ist den verschiedenen Definitionen die Vernetzung technischer Geräte und Systeme innerhalb eines privaten Wohnumfelds~\cite{Q2,Q3}.
Ein Smart Home beschreibt ein Wohngebäude oder eine Wohnung, in der Geräte und technische Systeme miteinander kommunizieren und automatisierte Funktionen bereitstellen.  
Zu diesen Systemen zählen beispielsweise Beleuchtung, Heizungssteuerung, Belüftung, Haushaltsgeräte, Unterhaltungselektronik sowie Kommunikationssysteme~\cite{Q1,Q2,Q3}.
Die Geräte orientieren sich an den Bedürfnissen der Bewohner:innen und ermöglichen eine automatisierte oder ferngesteuerte Bedienung.   
Dabei können Funktionen über Smartphones, Computer oder andere Netzwerkschnittstellen gesteuert und überwacht werden~\cite{Q1,Q3}.

Ein wesentliches Merkmal eines Smart Homes stellt die Vernetzung einzelner Komponenten dar.  
Durch den Datenaustausch zwischen den Geräten entstehen zusätzliche Funktionen und Dienste, die über den ursprünglichen Funktionsumfang einzelner Geräte hinausgehen.
Die Vernetzung ermöglicht beispielsweise die automatische Anpassung von Heizungs- oder Beleuchtungssystemen an bestimmte Zustände innerhalb des Gebäudes.  
Dadurch lassen sich Komfortfunktionen, Energieeinsparungen und Assistenzfunktionen realisieren.
Neben klassischen Anwendungen der Hausautomation gehören auch internetfähige Haushaltsgeräte und Systeme der Unterhaltungselektronik zum Smart Home.  
Dazu zählen unter anderem vernetzte Lautsprechersysteme, Smart-TV oder Haushaltsgeräte mit Zustandsüberwachung~\cite{Q1,Q2}.

Im Mittelpunkt moderner Smart-Home-Konzepte steht nicht ausschließlich eine zentrale Steuerung aller Komponenten.  
Vielmehr übernehmen einzelne Teilsysteme eigenständig bestimmte Aufgaben und tauschen die dafür notwendigen Informationen untereinander aus.  
Dieses Konzept wird als verteilte Intelligenz bezeichnet~\cite{Q2}. 
Viele Smart-Home-Geräte sind Bestandteil des \gls{iot}.  
Die Geräte besitzen Netzwerkfunktionen und kommunizieren über lokale Netzwerke oder das Internet miteinander~\cite{Q1,Q3}.
Durch die zunehmende Vernetzung entstehen neben funktionalen Vorteilen auch Anforderungen an Datenschutz und \gls{it}-Sicherheit.  
Da zahlreiche Geräte dauerhaft mit dem Internet verbunden sind, bestehen vergleichbare Risiken wie bei anderen internetfähigen Endgeräten~\cite{Q1}.

Zusätzlich unterscheiden sich Smart-Home-Systeme hinsichtlich ihrer technischen Umsetzung und ihrer Kommunikationsstandards.  
Aktuell dominieren häufig proprietäre Systeme und herstellerspezifische Insellösungen.  
Dadurch entstehen Einschränkungen bei der Kompatibilität unterschiedlicher Geräte und Systeme~\cite{Q2}.
Zusammenfassend definiert sich ein Smart Home durch die intelligente Vernetzung technischer Komponenten innerhalb eines privaten Wohnumfelds.  
Die Kombination aus Automatisierung, Kommunikation und digitaler Steuerung ermöglicht erweiterte Funktionen zur Steigerung von Komfort, Energieeffizienz und Bedienbarkeit~\cite{Q1,Q2,Q3}.